\chapter{KIỂM THỬ VÀ TRIỂN KHAI}

\section{Kiểm thử phần mềm}

\subsection{Chiến lược kiểm thử}
Dự án áp dụng phương pháp kiểm thử multi-layer:

\begin{enumerate}
    \item \textbf{Unit Tests:} Kiểm thử từng method riêng lẻ
    \item \textbf{Integration Tests:} Kiểm thử giao tiếp giữa các modules
    \item \textbf{API Tests:} Kiểm thử endpoints
    \item \textbf{UI/E2E Tests:} Kiểm thử toàn bộ flows (end-to-end)
\end{enumerate}

\subsection{Unit Tests}

\subsubsection{Framework}
Sử dụng \textbf{xUnit} kết hợp \textbf{Moq} library:

\begin{verbatim}
// PostsService.Tests.cs
public class PostsServiceTests
{
    private readonly Mock<IRepository<Post>> _mockRepository;
    private readonly PostsService _service;
    
    public PostsServiceTests()
    {
        _mockRepository = new Mock<IRepository<Post>>();
        _service = new PostsService(_mockRepository.Object);
    }
    
    [Fact]
    public async Task CreatePost_WithValidData_ReturnsPost()
    {
        // Arrange
        var request = new CreatePostRequest 
        { 
            Content = "Test post",
            UserId = "user123"
        };
        
        // Act
        var result = await _service.CreatePost(request);
        
        // Assert
        Assert.NotNull(result);
        Assert.Equal("Test post", result.Content);
        _mockRepository.Verify(
            x => x.AddAsync(It.IsAny<Post>()), 
            Times.Once);
    }
    
    [Fact]
    public async Task CreatePost_WithEmptyContent_ThrowsException()
    {
        // Arrange
        var request = new CreatePostRequest { Content = "" };
        
        // Act & Assert
        await Assert.ThrowsAsync<ArgumentException>(
            () => _service.CreatePost(request));
    }
}
\end{verbatim}

\subsubsection{Coverage Goals}
\begin{itemize}
    \item \textbf{Target Coverage:} Minimum 60\% cho services
    \item \textbf{Critical Services:} AuthService, PostsService - 80\%+
    \item \textbf{Test Methods:} Happy path + error cases + edge cases
    \item \textbf{Execution:} CI/CD pipeline chạy tests tự động
\end{itemize}

\subsection{Integration Tests}

\begin{verbatim}
// UsersService.Integration.Tests.cs
[Collection("Database collection")]
public class UsersServiceIntegrationTests
{
    private readonly IServiceProvider _serviceProvider;
    
    [Fact]
    public async Task UpdateUserProfile_SavesCorrectly()
    {
        // Arrange
        using (var scope = _serviceProvider.CreateScope())
        {
            var dbContext = scope.ServiceProvider
                .GetRequiredService<AppDbContext>();
            
            var user = new ApplicationUser 
            { 
                Id = "user1",
                DisplayName = "John Doe",
                Bio = "Old bio"
            };
            await dbContext.Users.AddAsync(user);
            await dbContext.SaveChangesAsync();
            
            var service = scope.ServiceProvider
                .GetRequiredService<IUsersService>();
            
            // Act
            await service.UpdateProfile("user1", new ProfileUpdateRequest
            {
                DisplayName = "Jane Doe",
                Bio = "New bio"
            });
            
            // Assert
            var updated = await dbContext.Users.FindAsync("user1");
            Assert.Equal("Jane Doe", updated.DisplayName);
            Assert.Equal("New bio", updated.Bio);
        }
    }
}
\end{verbatim}

\subsection{API Tests}

\begin{verbatim}
// PostsController.API.Tests.cs
public class PostsControllerTests
{
    private readonly HttpClient _client;
    
    [Fact]
    public async Task CreatePost_ReturnsCreated()
    {
        // Arrange
        var token = await GetJwtToken("user123");
        _client.DefaultRequestHeaders
            .Authorization = new AuthenticationHeaderValue(
                "Bearer", token);
        
        var request = new CreatePostRequest 
        { 
            Content = "My first post" 
        };
        var json = JsonSerializer.Serialize(request);
        var content = new StringContent(json, 
            Encoding.UTF8, "application/json");
        
        // Act
        var response = await _client.PostAsync("/api/posts", content);
        
        // Assert
        Assert.Equal(HttpStatusCode.Created, response.StatusCode);
        var responseContent = await response.Content.ReadAsStringAsync();
        var result = JsonSerializer
            .Deserialize<PostResponse>(responseContent);
        Assert.Equal("My first post", result.Content);
    }
    
    [Fact]
    public async Task GetPost_WithoutAuth_ReturnUnauthorized()
    {
        // Act
        var response = await _client.GetAsync("/api/posts/1");
        
        // Assert
        Assert.Equal(HttpStatusCode.Unauthorized, response.StatusCode);
    }
}
\end{verbatim}

\subsection{E2E/UI Tests}

Sử dụng Playwright hoặc Cypress cho browser automation:

\begin{verbatim}
// e2e/user-flow.spec.ts (Playwright)
test.describe('User Social Flow', () => {
    test('User can register, login, post, and like', async ({ page }) => {
        // Register
        await page.goto('http://localhost:3000/register');
        await page.fill('input[name=email]', 'test@example.com');
        await page.fill('input[name=password]', 'password123');
        await page.click('button:has-text("Register")');
        await page.waitForURL('http://localhost:3000/');
        
        // Login
        await page.goto('http://localhost:3000/login');
        await page.fill('input[name=email]', 'test@example.com');
        await page.fill('input[name=password]', 'password123');
        await page.click('button:has-text("Login")');
        await page.waitForURL('http://localhost:3000/home');
        
        // Create post
        await page.fill('textarea', 'My first post!');
        await page.click('button:has-text("Post")');
        await page.waitForSelector('.post-card');
        
        // Like post
        await page.click('button:has-text("❤️")');
        await page.waitForSelector('text=1 Likes');
    });
});
\end{verbatim}

\section{Triển khai CI/CD}

\subsection{GitHub Actions}

\begin{verbatim}
# .github/workflows/ci-cd.yml
name: CI/CD Pipeline

on:
  push:
    branches: [main, develop]
  pull_request:
    branches: [main]

jobs:
  build-backend:
    runs-on: ubuntu-latest
    steps:
      - uses: actions/checkout@v3
      
      - name: Setup .NET
        uses: actions/setup-dotnet@v3
        with:
          dotnet-version: 10.0.x
      
      - name: Restore
        run: dotnet restore backend/
      
      - name: Build
        run: dotnet build backend/ --no-restore
      
      - name: Run Tests
        run: dotnet test backend/ --no-build 
                    --verbosity normal 
                    --logger "trx;LogFileName=test-results.trx"
      
      - name: Upload Coverage
        uses: codecov/codecov-action@v3
        with:
          files: ./coverage/coverage.cobertura.xml
      
      - name: Publish
        run: dotnet publish backend/ -c Release 
                    -o ./publish
      
      - name: Build Docker Image
        run: |
          docker build -t ${{ secrets.REGISTRY_NAME }}/api:${{ 
            github.sha }} .
          docker push ${{ secrets.REGISTRY_NAME }}/api:${{ 
            github.sha }}

  build-frontend:
    runs-on: ubuntu-latest
    steps:
      - uses: actions/checkout@v3
      
      - name: Setup Node
        uses: actions/setup-node@v3
        with:
          node-version: 18.x
      
      - name: Install
        run: npm ci --prefix frontend
      
      - name: Lint
        run: npm run lint --prefix frontend
      
      - name: Build
        run: npm run build --prefix frontend
        env:
          REACT_APP_API_URL: ${{ secrets.API_URL }}
      
      - name: Upload to Azure Blob Storage
        run: |
          az storage blob upload-batch \
            -d web \
            -s frontend/dist \
            -c ${{ secrets.STORAGE_CONTAINER }}
      
      - name: Purge CDN
        run: az cdn endpoint purge -g ${{ secrets.RESOURCE_GROUP }} 
                     -n ${{ secrets.CDN_ENDPOINT }} 
                     --profile-name ${{ secrets.CDN_PROFILE }}

  deploy:
    needs: [build-backend, build-frontend]
    runs-on: ubuntu-latest
    if: github.ref == 'refs/heads/main'
    steps:
      - name: Deploy to Azure App Service
        uses: azure/webapps-deploy@v2
        with:
          app-name: ${{ secrets.APP_SERVICE_NAME }}
          publish-profile: ${{ secrets.AZURE_PUBLISH_PROFILE }}
          images: ${{ secrets.REGISTRY_NAME }}/api:${{ github.sha }}
\end{verbatim}

\subsection{Build Stages}

\begin{enumerate}
    \item \textbf{Source Stage:} Pull code từ GitHub
    \item \textbf{Build Stage:} Compile \& run tests
    \item \textbf{Quality Gate:} Verify test coverage \& code quality
    \item \textbf{Package Stage:} Create Docker images
    \item \textbf{Deploy Stage:} Deploy to Azure (if main branch)
\end{enumerate}

\section{Triển khai Cloud}

\subsection{Azure Infrastructure}

\begin{verbatim}
┌─────────────────────────────────────┐
│       Azure Container Registry      │
│    (Docker Images Storage)          │
└──────────────┬──────────────────────┘
               │
┌──────────────▼──────────────────────┐
│    Azure App Service                │
│  (Web API + Background Jobs)        │
│  - Pricing Tier: B2 or higher       │
│  - Instances: 2+ for HA             │
│  - Auto-scale: 3-10 instances       │
└──────────────┬──────────────────────┘
               │
    ┌──────────┼──────────┐
    │          │          │
┌───▼──┐  ┌────▼────┐  ┌─▼────────────┐
│ DNS  │  │  Azure  │  │ Application  │
│  CDN │  │   SQL   │  │  Insights    │
│      │  │Database │  │  Monitoring  │
└──────┘  │         │  └──────────────┘
          │- Premium│
          │ - Geo- │
          │ Replicat
          │ed       │
          └─────────┘
\end{verbatim}

\subsection{Azure Services Used}

\begin{enumerate}
    \item \textbf{Azure App Service:}
    \begin{itemize}
        \item Host cho ASP.NET Core API
        \item Auto-scaling based on CPU/Memory
        \item Built-in deployment slots (staging)
        \item Application Insights integration
    \end{itemize}

    \item \textbf{Azure SQL Database:}
    \begin{itemize}
        \item Managed SQL Server database
        \item Automatic backups \& geo-replication
        \item Point-in-time restore
        \item Firewall rules \& VNet integration
    \end{itemize}

    \item \textbf{Azure Blob Storage:}
    \begin{itemize}
        \item Store user uploaded images
        \item CDN integration for fast delivery
        \item Lifecycle management (auto-delete old files)
        \item Access via SAS tokens (secure URLs)
    \end{itemize}

    \item \textbf{Azure Key Vault:}
    \begin{itemize}
        \item Store connection strings \& secrets
        \item Rotate keys automatically
        \item Audit access logs
        \item Access via managed identity
    \end{itemize}

    \item \textbf{Azure Application Insights:}
    \begin{itemize}
        \item Application Performance Monitoring (APM)
        \item Track requests, dependencies, exceptions
        \item Custom metrics \& events
        \item Alerting based on thresholds
    \end{itemize}

    \item \textbf{Azure CDN:}
    \begin{itemize}
        \item Cache static assets (React bundles)
        \item Cache API responses
        \item Edge server locations worldwide
        \item DDoS protection
    \end{itemize}
\end{enumerate}

\subsection{Deployment Process}

\subsubsection{Step 1: Prepare Infrastructure}
\begin{verbatim}
# Using Azure Resource Manager template (Bicep)
az deployment group create \
  --resource-group my-rg \
  --template-file main.bicep \
  --parameters environment=prod
\end{verbatim}

\subsubsection{Step 2: Deploy Application}
\begin{verbatim}
# Backend deployment
az webapp deployment container config \
  --name interacthub-api \
  --resource-group my-rg \
  --docker-registry-server-url registry.url \
  --docker-registry-server-username username \
  --docker-registry-server-password password

az webapp deployment container config \
  --name interacthub-api \
  --settings DOCKER_REGISTRY_SERVER_URL=...
\end{verbatim}

\subsubsection{Step 3: Configure Environment}
\begin{verbatim}
# Set connection strings via Key Vault
az webapp config appsettings set \
  --name interacthub-api \
  --resource-group my-rg \
  --settings \
    ConnectionStrings__DefaultConnection=@Microsoft.KeyVault(\
      SecretsUri=https://myvault.vault.azure.net/\
      secrets/db-conn-string/) \
    Jwt__Key=@Microsoft.KeyVault(\
      SecretsUri=https://myvault.vault.azure.net/\
      secrets/jwt-key/)
\end{verbatim}

\subsubsection{Step 4: Run Migrations}
\begin{verbatim}
# Connect to Azure App Service Kudu
# And run migrations
dotnet ef database update \
  -c AppDbContext \
  --connection "connection-string"
\end{verbatim}

\section{Monitoring \& Troubleshooting}

\subsection{Application Insights}

\begin{itemize}
    \item \textbf{Request Duration:} Monitor API response times
    \item \textbf{Error Rate:} Track exceptions \& failed requests
    \item \textbf{Dependencies:} Database, external API calls
    \item \textbf{Performance Counters:} CPU, Memory, Disk I/O
    \item \textbf{Custom Metrics:} Business metrics (active users, posts/day)
\end{itemize}

\subsection{Alerts}

\begin{verbatim}
// Example alerts configured
- If error rate > 5% for 5 minutes → Alert
- If response time > 2 seconds (p95) → Alert
- If CPU > 80% for 10 minutes → Alert
- If database connections > 90 threshold → Alert
\end{verbatim}

\subsection{Scaling Strategy}

\begin{itemize}
    \item \textbf{Vertical Scaling:} Upgrade App Service tier (B2 → P1V2)
    \item \textbf{Horizontal Scaling:} Auto-scale 3-10 instances
    \item \textbf{Database Scaling:} Upgrade SQL Database tier
    \item \textbf{Cache:} Add Redis for frequently accessed data
    \item \textbf{CDN:} Cache static assets at edge locations
\end{itemize}
